\begingroup\small\setlength{\tabcolsep}{3pt}\def\LTcaptype{quadro}
\setlength{\LTleft}{0pt}
\setlength{\LTcapwidth}{\textwidth}\renewcommand{\arraystretch}{1.12}
\begin{longtable}{|>{\raggedright\arraybackslash}p{2.8cm}|>{\raggedright\arraybackslash}p{3.3cm}|>{\raggedright\arraybackslash}p{4.1cm}|>{\raggedright\arraybackslash}p{4.5cm}|}
\caption{Finalidade, técnica e população dos estudos incluídos.}\label{qua:sintese-tipologia}\\\hline
\textbf{Estudo} & \textbf{Finalidade principal} & \textbf{Família de técnica} & \textbf{Nível e n (unidade)} \\ \hline
\endfirsthead
\multicolumn{4}{c}{\quadroname\ \thequadro\ -- Continuação}\\\hline
\textbf{Estudo} & \textbf{Finalidade principal} & \textbf{Família de técnica} & \textbf{Nível e n (unidade)} \\ \hline
\endhead
\hline\endfoot\hline\endlastfoot
\citeonline{Math2021_001} & Predição/\allowbreak{}classificação & Rede neural \textit{feedforward}/\allowbreak{}MLP e \textit{Random Forest}; seleção manual/\allowbreak{}RF/\allowbreak{}qui-quadrado/\allowbreak{}eliminação recursiva; SMOTE no treino. & Escolar, avaliação internacional PISA 2012 computadorizada.; 271.323 registros CBA iniciais; 58.753 estudantes no conjunto analítico. \\ \hline
\citeonline{Implementation2025_000} & Predição/\allowbreak{}classificação & Dez classificadores (árvores, discriminante, SVM, regressão logística, KNN, redes e combinações de modelos); combinação SMOTE/\allowbreak{}RUSBoost destacada. & Ensino primário, escola privada; atributos da quarta série e nota de conclusão da sexta.; 280 estudantes; 33 atributos originais, 48 colunas após codificação. \\ \hline
\citeonline{Multimodels2020_002} & Predição/\allowbreak{}classificação & EDM, ML e métodos estatísticos, conforme registro histórico; algoritmos específicos não reconferidos. & Ensino médio, escolas do Camboja, conforme metadados/\allowbreak{}histórico.; Não reconferido; fonte primária indisponível. \\ \hline
\citeonline{Analysis2022_003} & Predição/\allowbreak{}classificação & DT, RF, \textit{Naive Bayes}, MLP e JRip; seleção \textit{wrapper} de atributos. & Ensino secundário, duas escolas portuguesas; Matemática e Português separados.; Dimensão amostral não informada suficientemente; não importar n da UCI por suposição. \\ \hline
\citeonline{Design2025_004} & Personalização/\allowbreak{}STI & Grafo de conhecimento; perfis por redes profundas/\allowbreak{}RNN; aprendizagem por reforço e análise de sentimento. & Não informado suficientemente; cursos de Matemática, Chinês e Inglês.; Não informado; amostragem, alocação e origem humana/\allowbreak{}simulada dos resultados incertas. \\ \hline
\citeonline{Identifying2017_006} & Predição/\allowbreak{}classificação & RF, J48, \textit{Naive Bayes}, MLP e regressão logística. & Oitavo ano, Turquia, TIMSS 2011; atributo idade 13-16.; 6.928 originais; 678 excluídos por dados ausentes/\allowbreak{}inexatos; 6.250 analíticos. \\ \hline
\citeonline{Innovative2023_005} & Personalização/\allowbreak{}STI & Grafo/\allowbreak{}perfil; ACO e CNN para trajetórias; TransR e CRF semissupervisionado/\allowbreak{}KNN para recursos. & Primeiro ano universitário, duas turmas de Gestão Financeira, Matemática Superior.; 120 iniciais; 50 experimentais e 52 controles válidos; 18 casos sem explicação. \\ \hline
\citeonline{Computational2017_008} & Modelagem cognitiva & Arquitetura Apprentice Learner; modelos DECISION TREE e TRESTLE; simulações e modelos mistos binomiais. & Frações e equações em dados educacionais secundários; tarefas de caixas e setas com participantes recrutados via Mechanical Turk. Níveis heterogêneos.; Frações: 79; equações: 71 estudantes, 10.052 passos disponíveis e 4.568 analisados; caixas e setas: 202 participantes. Não somar bases nem frações reutilizadas. \\ \hline
\citeonline{Machine2024_009} & Predição/\allowbreak{}classificação & Árvore de decisão podada, MATLAB. & Primeiro ano de Engenharia Industrial, Universidade de East Sarajevo.; 50 estudantes; dados de 2017-2021. \\ \hline
\citeonline{Enhancing2025_012} & Intervenção didática & Animação computacional e aprendizagem cooperativa \textit{jigsaw}; não experimento de ML. & Segundo ano senior high school/\allowbreak{}Form2, Matemática eletiva, duas escolas públicas de Ghana.; 196; 95 \textit{jigsaw}/\allowbreak{}animação, 101 tradicional; uma escola por condição. \\ \hline
\citeonline{He2025_6915} & Intervenção didática & Escolar: GeoGebra e duração orientada por modelo. Outros componentes: meta-análise e regressão/\allowbreak{}EDM, incluindo XGBoost. & Escolar: ensino secundário/\allowbreak{}geometria analítica, escola no sudeste da China; literatura: K-12.; Literatura: 423 estudos/\allowbreak{}efeitos, não participantes ou novas inclusões do TCC. Escolar: oito turmas de 45, mas total declarado90; não reconciliado. \\ \hline
\citeonline{Villegas2025_6916} & Personalização/\allowbreak{}STI & STI; \textit{deep learning}/\allowbreak{}NLP, \textit{transformers}/\allowbreak{}RNN/\allowbreak{}BERT; modelos supervisionados DNN/\allowbreak{}SVM descritos; \textit{feedback}/\allowbreak{}sequência adaptativos. & Universitário STEM: Matemática, Física e Programação.; 450 no conjunto STEM; n específico de Matemática/\allowbreak{}grupos não estabelecido nas passagens. Entrevistas com estudantes/\allowbreak{}educadores, n não suficientemente localizado. \\ \hline
\citeonline{Ozseven2026_6917} & Predição/\allowbreak{}classificação & ANFIS/\allowbreak{}neuro-fuzzy, RF e XGBoost; saídas contínuas, híbrido ponderado, \textit{stacking}/\allowbreak{}regressão logística e limiares binários. & ASSISTments descrito para middle/\allowbreak{}high school; recorte analítico não suficientemente delimitado.; n final de estudantes/\allowbreak{}interações não localizado; filtros menos de5 e limite1000, perdas não quantificadas. \\ \hline
\citeonline{UniversityMathematics2026_6919} & Integração curricular & Fourier aplicado à metrologia SiC; questionário adaptativo denominado IA; PCA para estrutura do questionário. Algoritmo ML não identificado nas passagens. & Segundo ano universitário, Measurement and Control Technology/\allowbreak{}Advanced Mathematics.; Controle2021 n94; pós-reforma2025 n89; 12 entrevistados do experimental, sem somá-los como novos participantes. \\ \hline
\citeonline{Echeveria2025_6920} & Predição/\allowbreak{}classificação & K-means para categorias; Boruta; SVM/\allowbreak{}NB/\allowbreak{}MLP; entrevistas docentes. & Junior high school, uma escola nas Filipinas.; 633 recrutados,428 completos,5 docentes; métodos2023-2024 versus resultados2022-2023 não reconciliados. \\ \hline
\end{longtable}\endgroup
\fonte{Elaborado pelo autor (2026).}
\begingroup\small\setlength{\tabcolsep}{3pt}\def\LTcaptype{quadro}
\setlength{\LTleft}{0pt}
\setlength{\LTcapwidth}{\textwidth}\renewcommand{\arraystretch}{1.12}
\begin{longtable}{|>{\raggedright\arraybackslash}p{2.8cm}|>{\raggedright\arraybackslash}p{3.3cm}|>{\raggedright\arraybackslash}p{4.1cm}|>{\raggedright\arraybackslash}p{4.5cm}|}
\caption{Alvos, validação e métricas dos estudos incluídos.}\label{qua:sintese-validacao}\\\hline
\textbf{Estudo} & \textbf{Alvo} & \textbf{Validação e fonte} & \textbf{Métrica ou resultado reportado} \\ \hline
\endfirsthead
\multicolumn{4}{c}{\quadroname\ \thequadro\ -- Continuação}\\\hline
\textbf{Estudo} & \textbf{Alvo} & \textbf{Validação e fonte} & \textbf{Métrica ou resultado reportado} \\ \hline
\endhead
\hline\endfoot\hline\endlastfoot
\citeonline{Math2021_001} & Três classes operacionais derivadas da razão pontos obtidos/\allowbreak{}máximos nos itens matemáticos; não são níveis oficiais do PISA ou escala individual validada. & Partição 70/\allowbreak{}30; seleção de atributos declarada só no treino; subdivisão treino/\allowbreak{}validação da RN. Outra passagem afirma monitorar teste no parada antecipada, impedindo confirmar independência. (páginas físicas 2–5, Tabelas IV–VI) & RN/\allowbreak{}eliminação recursiva: acurácia denominada \textit{macro-average} pelos autores 78,39\% e F1 0,68 (Tabela IV). Kappa máximo 0,46 (Tabela VI); agregações distintas. \\ \hline
\citeonline{Implementation2025_000} & Low/\allowbreak{}Medium/\allowbreak{}High da nota matemática da sexta série; dados da quarta em 2017/\allowbreak{}2018 e desfecho em 2020/\allowbreak{}2021. & Partição 70/\allowbreak{}30; validação em cinco partes no treino; reamostragem declarada no treino. (páginas físicas 2–3 e 7–8, Tabelas 8–9) & SMOTE+Ensemble RUSBoosted Trees: teste 75\%; classe Low recall 43\%, precisão 67\%, F-measure 52\% (Tabela 8). Validação 92\% é outra etapa. \\ \hline
\citeonline{Multimodels2020_002} & Desempenho matemático conforme título; operacionalização não reconferida. & Não reconferida. (registro histórico; fonte primária não reconferida) & Não reconferida; nenhum valor ratificado nesta versão. \\ \hline
\citeonline{Analysis2022_003} & Nota final em aprovação/\allowbreak{}reprovação ou cinco classes. & Protocolo próprio de partição/\allowbreak{}validação não especificado suficientemente. (páginas físicas 2–5, Tabelas 6 e 8) & Matemática: Tabela 6, DT multiclasses 85,62\% antes e 88,75\% após \textit{wrapper}; Tabela 8, RF binária 93,49\% antes e 94,10\% após. \\ \hline
\citeonline{Design2025_004} & Selecionar próximo nó da trajetória segundo conhecimento, progresso e emoção; ganhos, tempo e satisfação. & Comparação com filtragem colaborativa e rede rasa; sem partição, alocação, análise inferencial ou incerteza suficientes. (páginas físicas 1–5, Tabela I) & Ganhos relativos reportados: aprendizagem 15\%, tempo 20\%; satisfação 4,2/\allowbreak{}5. Sem incerteza suficiente. \\ \hline
\citeonline{Identifying2017_006} & MATHACH binário: primeira valor plausível acima de 500 ou não. & Dez dobras declaradas; texto diz uma dobra treino/\allowbreak{}nove teste. Não determina se erro de relato ou execução. (páginas físicas 4–5, 9 e 11) & Regressão logística: kappa 0,4885; RMSE 0,3885 (Tabela 3); F-measure ponderada 0,779 e ROC-AUC 0,817 (Tabela 4). \\ \hline
\citeonline{Innovative2023_005} & Trajetórias/\allowbreak{}recursos personalizados; pré/\allowbreak{}pós de Matemática e percepções. & Comparação pré/\allowbreak{}pós entre turmas; alocação não localizada; p basal 0,583 não prova equivalência. (páginas físicas 9–12, Tabelas 1–6) & Experimental 63,795 para 75,631; controle 62,865 para 66,314. Narrativa usa também 74,631; pré-teste com gl=80 não reconciliado com 50+52. \\ \hline
\citeonline{Computational2017_008} & Comportamento/\allowbreak{}curvas humanos e efeitos qualitativos de sequenciamento; ajuste de mecanismos computacionais. & Agentes recebem problemas na ordem dos registros; curvas/\allowbreak{}resíduos e ajuste. Não é nova intervenção humana prospectiva; sete domínios não são todos matemáticos. (páginas impressas 44–48 e 52–56) & Frações: AIC 9.512,8 (TRESTLE) e 9.521,3 (árvore); ajuste de curvas e predições instrucionais qualitativas, não escala de proficiência. \\ \hline
\citeonline{Machine2024_009} & Resultado matemático anual fail/\allowbreak{}conditional/\allowbreak{}pass; preditores incluem notas de outras disciplinas e atividades. & Não localizado protocolo próprio treino/\allowbreak{}teste ou validação cruzada; passagem k-fold descreve outro estudo. (páginas físicas 4–7) & Acurácia reportada 72\%, erro 28\%; não apresentar como teste independente. \\ \hline
\citeonline{Enhancing2025_012} & Pré/\allowbreak{}pós em teoremas de círculo. & Pré/\allowbreak{}pós e Mann-Whitney; §2.4 declara sorteio de escolas, §2.7 nega alocação aleatória. Comparações separadas não ajustam medida inicial/\allowbreak{}escola. (páginas físicas 10–15, Tabela 3) & Pré-teste: U=6.714, p inferior a 0,001, efeito 0,346; pós: U=7.948,5, p inferior a 0,001, efeito 0,568. Comparações não ajustadas. \\ \hline
\citeonline{He2025_6915} & Escolar: resultados após 18 períodos com GeoGebra. Modelagem: tamanho de efeito de tecnologia educacional, não proficiência individual. & Dois controles/\allowbreak{}seis experimentais, testes t pré/\allowbreak{}pós; seis efeitos escolares revalidam modelo. Professor confundido com condição; agrupamento por turma não esclarecido. (páginas físicas 11–16 e 31–35, Tabelas 11–13) & Aplicação: Hedges g entre 0,727 e 0,764 nos seis grupos (Tabela 13). Não contar como seis estudos novos. \\ \hline
\citeonline{Villegas2025_6916} & Feedback, sequência, precisão, progresso, sucesso em problemas, tempo e percepções; construtos distintos. & Experimental/\allowbreak{}controle, t/\allowbreak{}ANOVA e entrevistas temáticas; recrutamento/\allowbreak{}alocação/\allowbreak{}completude insuficientes. Semestre na introdução, seis semanas métodos, semana8 somente Programação. (páginas físicas 2, 13–16 e 19–26) & Tabela4 Matemática: precisão do feedback80,2\% para88,7\%, melhoria relativa10,6\%. Tabela3 global STEM88,5\% versus76,3\%; não atribuir somente a Matemática. \\ \hline
\citeonline{Ozseven2026_6917} & correct de cada interação estudante-problema; não proficiência global ou classe validada de risco. & 60/\allowbreak{}20/\allowbreak{}20 por estudante; cada estudante em único subconjunto declarado. Disponibilidade temporal dos atributos e validade dos construtos exigem esclarecimento. (páginas físicas 3–4, 7 e 10–12) & ANFIS: acurácia 0,922063, revocação de acerto 0,9962. Híbrido ponderado: RMSE 0,2414. Tarefas e agregações distintas. \\ \hline
\citeonline{UniversityMathematics2026_6919} & Integração matemática/\allowbreak{}indústria/\allowbreak{}sustentabilidade; experiência de aula, apoio, ativação cognitiva e estrutura disciplinar; entrevistas. & Coortes não contemporâneas; testes t; alfa/\allowbreak{}KMO/\allowbreak{}Bartlett/\allowbreak{}PCA do questionário não validam mecanismo de personalização. Não significância basal não prova equivalência. (páginas físicas 7–8, 12–14 e 18–19; leitura parcial) & Questionário: tamanhos de efeito de -0,38 a 0,67 nas quatro dimensões. Alfa 0,906; KMO 0,643. Não mede efeito isolado da IA. \\ \hline
\citeonline{Echeveria2025_6920} & Quatro categorias de mentalidade derivadas de questionário; não domínio matemático ou resposta à intervenção. & Relato menciona dois terços,80/\allowbreak{}20 e dez dobras; execução final incerta. Oito itens do alvo permanecem nos preditores. (páginas físicas 3 e 6–12) & SVM: 83/\allowbreak{}84=98,8\%, Kappa 0,98. Matriz MLP tem total impresso 84 e colunas somando 86; não harmonizada. \\ \hline
\end{longtable}\endgroup
\fonte{Elaborado pelo autor (2026).}
\begin{table}[htbp]\centering
\caption{Frequência da finalidade principal no corpus de 15 estudos.}\label{tab:tipologia-frequencias}
\begin{tabular}{lr}\hline Finalidade & Estudos (n) \\ \hline
Predição/\allowbreak{}classificação & 8 \\
Personalização/\allowbreak{}STI & 3 \\
Modelagem cognitiva & 1 \\
Intervenção didática & 2 \\
Integração curricular & 1 \\
\hline Total & 15 \\ \hline\end{tabular}
\fonte{Elaborado pelo autor (2026).}\end{table}
